% ===================================================================
%  TABLICA Nr 11 — Miasto i jego zabytki. — Pożar.
%  Meme regime que les precedents. Premiere de la TROISIEME SERIE :
%  l'apparat porte « SERIA TRZECIA », ordinal apres le nom, comme aux
%  tableaux 1 et 7 — c'est le membre polonais de SERIO qui le lit.
%
%  DEUX BLOCS PORTENT DEUX ALINEAS SOUS UNE SEULE CLE, t11-c1-12-1 et
%  t11-c2-03-1 : le fac-simile ido y a une coupure de page, pas un
%  paragraphe. Les traductions les soudent, le neerlandais comme le
%  tcheque, et ordo.py connait le cas.
%
%  LE 9 EST LE PLUS GRAS DU LIVRET ET LE PLUS ETRANGE : le
%  compositeur y met en gras des mots-outils entiers — « quale »,
%  « la sistemo di traciono di la », « La », « tirata », « urbo al
%  altra ». Aucun n'est un terme de planche. On les reproduit parce
%  que toutes les autres colonnes les reproduisent : le gras du
%  fac-simile fait partie du document.
%
%  ET LE 6 DE LA DEUXIEME SCENE DONNE DEUX FOIS LE RENVOI (81), aux
%  musiciens puis au violon. Ce n'est pas une faute a corriger : la
%  planche numerote ainsi, et les colonnes qui s'en ecarteraient
%  casseraient le releve.
% ===================================================================

%%K t11-apar-1 apar
\VUsaut{2.0mm}
\VUcentre{13.2pt}{90}{SERIA TRZECIA}
\VUsaut{3.0mm}
\VUfilet{16mm}
\VUsaut{5.6mm}
\VUtitre{15.0pt}{60}{TABLICA Nr 11}
\VUsaut{1.4mm}
\VUpk{11.4pt}{40}{Miasto i jego zabytki. --- Pożar.}
\VUsaut{2.2mm}

%%K t11-tit-1 sub
\VUpk{10.2pt}{40}{Przedmieście.}
\VUsaut{3.0mm}

%%K t11-c1-01-1 p
1. --- Mieszkam w wielkim mieście położonym o 100 kilometrów od Oceanu, a które mimo to
jest \VUgras{portem} \textsuperscript{(1)} pierwszej rangi. Rzeka, która je obrzeża, ma
szerokość 400 metrów i tworzy bardzo piękną redę.

%%K t11-c1-02-1 p
2. --- Cała część miasta, która leży przy \VUgras{nabrzeżach} \textsuperscript{(2)},
wygląda całkiem morsko i przemysłowo i tworzy \VUgras{przedmieście}
\textsuperscript{(3)}, zamieszkane tylko przez żeglarzy i robotników. Ci pracują w
\VUgras{fabrykach} \textsuperscript{(4)} i w \VUgras{gazowni} \textsuperscript{(5)},
której wysoki \VUgras{komin} \textsuperscript{(6)} i \VUgras{gazometr} widać z bardzo
daleka.

%%K t11-c1-03-1 p
3. --- W średniowieczu miasto było ufortyfikowane i znajduje się jeszcze kilka
pozostałości jego wałów, zwłaszcza \VUgras{brama} \textsuperscript{(8)} starego
\VUgras{zamku warownego} \textsuperscript{(7)}, obrzeżona swoimi dwiema \VUgras{wieżami}
\textsuperscript{(9)}, które teraz służą za \VUgras{więzienie}.

%%K t11-c1-04-1 p
4. --- W całej tej \VUgras{dzielnicy}, która wygląda bardzo staro, tylko jedna budowla
jest względnie nowoczesna: to \VUgras{komora celna} \textsuperscript{(10)}. Leży ona
między nabrzeżem a \VUgras{uliczką} \textsuperscript{(11)}, która prowadzi do starego
\VUgras{ratusza} \textsuperscript{(12)}. Łatwo poznaje się ten ostatni zabytek po
olbrzymiej \VUgras{dzwonnicy} \textsuperscript{(13)}, która panuje nad starym miastem.

%%K t11-c1-05-1 p
5. --- Po drugiej stronie ulicy \VUgras{stoi} budynek zbudowany w tym samym stylu co
komora celna: to \VUgras{Giełda} \textsuperscript{(14)}. Jedna z jej fasad patrzy na mały
plac, gdzie jest \VUgras{prefektura} \textsuperscript{(15)}. Nasze miasto bowiem, nie
będąc stolicą, jest mimo to ważną stolicą departamentu.

%%K t11-tit-2 sub
\VUsaut{5.0mm}
\VUpk{10.2pt}{40}{Najwyższa część miasta.}
\VUsaut{3.4mm}

%%K t11-c1-06-1 p
6. --- Garnizon, dość liczny, jest zakwaterowany w rozległych, bardzo zdrowych
\VUgras{koszarach} \textsuperscript{(16)}; rozciągają się one aż do krat \VUgras{parku
miejskiego} \textsuperscript{(17)}. \VUgras{Bulwar} \textsuperscript{(19)}, który
prowadzi do tego parku, biegnie za \VUgras{kasą oszczędnościową} \textsuperscript{(18)} i
dochodzi do placu \VUgras{katedry} \textsuperscript{(20)}.

%%K t11-c1-07-1 p
7. --- Wszystkie te zabytki rozciągają się amfiteatralnie na wzgórku, gdzie jest też
nasze rozległe \VUgras{liceum} \textsuperscript{(21)}.

%%K t11-c1-07-2 p
Jeśli schodzi się nieco pochyłą drogą, która dołącza do Wielkiej ulicy, dochodzi się do
\VUgras{Sądu} \textsuperscript{(22)}, przed którym rozciąga się przyjemny \VUgras{ogród
publiczny} \textsuperscript{(26)}. Ten \VUgras{skwer} nie jest bardzo duży, ale tworzy w
środku miasta oazę zieleni i chłodu. Dlatego jest bardzo uczęszczany przez mieszczan,
którzy lubią przychodzić i siadać pod namiotem \VUgras{kawiarni}
\textsuperscript{(25)}, żeby słuchać muzyków wojskowych, którzy w czwartki i niedziele
grają w \VUgras{kiosku} \textsuperscript{(27)} postawionym między trawnikami a kępami
krzewów.

%%K t11-c1-08-1 p
8. --- Na jednym z tych trawników stoi brązowy \VUgras{posąg} \textsuperscript{(28)},
pomnik wzniesiony ku pamięci jednego z naszych współmieszczan, który oddał wielkie
usługi swojemu rodzinnemu miastu.

%%K t11-tit-3 sub
\VUsaut{4.0mm}
\VUpk{10.2pt}{40}{Środki transportu.}
\VUsaut{3.4mm}

%%K t11-c1-09-1 p
9. --- Dotąd nasze miasto nie jest jeszcze oświetlane światłem elektrycznym, ma tylko
\VUgras{latarnie gazowe} \textsuperscript{(29)}. Ale \VUgras{tutejsza prasa} prowadzi
kampanię, żeby ten system oświetlenia został ulepszony, \VUgras{jak} już udoskonalono
\VUgras{system trakcyjny} tramwajów. \VUgras{Dawne} wozy, \VUgras{ciągnione} przez
\VUgras{konie,} zostały praktycznie zastąpione przez \VUgras{tramwaje elektryczne}
\textsuperscript{(31)}, które szybko przewożą podróżnych z jednego krańca \VUgras{miasta
na drugi,} za bardzo \VUgras{tanią} cenę.

%%K t11-c1-10-1 p
10. --- Blisko Sądu jest \VUgras{Bank} \textsuperscript{(32)}, gdzie dyskontuje się walory
handlowe. \VUgras{Urzędnicy bankowi} \textsuperscript{(33)} chodzą inkasować należności w
domach.

%%K t11-tit-4 sub
\VUsaut{4.0mm}
\VUpk{10.2pt}{40}{Sztuka i nauczanie.}
\VUsaut{3.4mm}

%%K t11-c1-11-1 p
11. --- Piękny budynek, który stoi obok, to \VUgras{muzeum.} Mieści ono zarazem
\VUgras{bibliotekę} \textsuperscript{(34)} miasta, piękne \VUgras{zbiory
przyrodnicze} \textsuperscript{(35)} (Muzeum przyrody) i wdzięczne \VUgras{muzeum
malarstwa} \textsuperscript{(36)}, które stanowi wspaniałą stałą \VUgras{wystawę}.

%%K t11-c1-11-2 p
Otwiera się ono dla publiczności dwa razy w tygodniu, ale cudzoziemcy mogą je zwiedzać
każdego dnia za napiwek dla \VUgras{strażnika} \textsuperscript{(37)}.
\VUgras{Malarze} \textsuperscript{(38)} przychodzą tam pracować. Widzi się ich
\VUgras{przed} ich sztalugą, z \VUgras{paletą} w ręku, kopiujących sławne obrazy.

%%K t11-c1-12-1 p
12. --- Na wzniesieniu, za muzeum, zaczyna się widzieć \VUgras{kopułę}
\textsuperscript{(40)} \VUgras{szpitala} \textsuperscript{(39)} i budowle
\VUgras{seminarium nauczycielskiego} \textsuperscript{(41)}, w którym kształcono
nauczycieli naszych szkół gminnych. Na placu Teatralnym jest \VUgras{urząd pocztowy}
\textsuperscript{(43)} umieszczony w \VUgras{domu prywatnym} \textsuperscript{(42)}.

%%K t11-tit-5 sub
\VUsaut{5.0mm}
\VUpk{11.4pt}{40}{Pożar.}
\VUsaut{3.0mm}

%%K t11-tit-6 sub
\VUpk{10.2pt}{40}{Okrzyk o pożarze.}
\VUsaut{3.4mm}

%%K t11-c2-01-1 p
1. --- Przerażająca wieść rozchodzi się po mieście: pożar właśnie wybuchł w teatrze! W
chwili, gdy przybywam na \VUgras{plac} \textsuperscript{(96)}, tłum jest już zebrany na
\VUgras{chodniku} \textsuperscript{(46)} i na \VUgras{jezdni} \textsuperscript{(47)}.
\VUgras{Urzędnicy pocztowi} \textsuperscript{(44)} i \VUgras{listonosze}
\textsuperscript{(45)} wychodzą z urzędu pocztowego. \VUgras{Motorniczy}
\textsuperscript{(48)} musiał zatrzymać swój wóz tramwajowy, żeby przepuścić miejską
\VUgras{karetkę} \textsuperscript{(50)}, która przybywa z \VUgras{chirurgiem}
\textsuperscript{(52)} i \VUgras{pielęgniarzem} \textsuperscript{(51)}, żeby opatrzyć
\VUgras{rannego} \textsuperscript{(53)} przyniesionego na \VUgras{noszach}
\textsuperscript{(54)} blisko \VUgras{kiosku z gazetami} \textsuperscript{(55)}.

%%K t11-c2-02-1 p
2. --- Kompania piechoty, która wracała z ćwiczeń, zatrzymała się, żeby zapewnić służbę
porządkową razem z \VUgras{konną policją miejską} \textsuperscript{(61)}.
\VUgras{Jeźdźcy,} których konie stają dęba, lepiej niż \VUgras{żołnierze}
\textsuperscript{(57)} czy \VUgras{żandarmi} \textsuperscript{(63)} cofają
\VUgras{ciekawskich} \textsuperscript{(56)}. Zresztą żandarmi są bardzo zajęci. Jeden z
nich wskazuje \VUgras{policjantom} \textsuperscript{(64)}, dokąd trzeba przewieźć innego
\VUgras{poszkodowanego} \textsuperscript{(65)}, dzielnego strażaka, który spadł z drabiny.

%%K t11-c2-03-1 p
3. --- \VUgras{Oficer} \textsuperscript{(66)} określa dokładnie miejsce, gdzie trzeba
ustawić przybywającą \VUgras{sikawkę parową} \textsuperscript{(67)}. Czekając na tę
potężną maszynę, ustawił w baterii \VUgras{sikawkę ręczną} \textsuperscript{(68)},
której długi \VUgras{wąż} \textsuperscript{(69)} rzuca wodę potokiem na ogień. Sześciu
strażaków nią manewruje, podczas gdy ich towarzysz, który trzyma \VUgras{prądownicę}
\textsuperscript{(70)}, kieruje \VUgras{strumień} \textsuperscript{(71)} ku środkowi
pożaru.

%%K t11-c2-04-1 p
4. --- Po drugiej stronie teatru oparto wielką \VUgras{drabinę} \textsuperscript{(72)}.
Pozwala ona strażakom zbliżyć się do płomieni, które tryskają wśród kłębów czarnego
\VUgras{dymu} \textsuperscript{(75)}.

%%K t11-tit-7 sub
\VUsaut{5.0mm}
\VUpk{10.2pt}{40}{Dzielne czyny.}
\VUsaut{3.4mm}

%%K t11-c2-05-1 p
5. --- \VUgras{Pożar} \textsuperscript{(74)} narodził się w foyer \VUgras{teatru}
\textsuperscript{(73)}, podczas gdy próbowano \VUgras{operę}, której pierwsze
\VUgras{przedstawienie} miało się odbyć jutro. Szczęśliwie tylko niewielu
\VUgras{widzów} \textsuperscript{(76)} było na widowni. Biedacy schronili się na
\VUgras{balkonie} \textsuperscript{(77)}, gdzie dzielni \VUgras{strażacy}
\textsuperscript{(86)} przychodzą im na ratunek z drabiną i linami.

%%K t11-c2-06-1 p
6. --- Pod \VUgras{perystylem} \textsuperscript{(78)}, gdzie \VUgras{afisz}
\textsuperscript{(79)} informuje o przedstawieniu, widzi się \VUgras{muzyków}
\textsuperscript{(81)} z \VUgras{orkiestry} \textsuperscript{(80)}, jak uciekają, niosąc
swoje instrumenty: \VUgras{skrzypce} \textsuperscript{(81)}, \VUgras{flet}
\textsuperscript{(82)}, \VUgras{wiolonczelę} \textsuperscript{(83)}, \VUgras{puzon}
\textsuperscript{(84)}, \VUgras{bęben} \textsuperscript{(85)}, itd.. Usiłuje się
uratować część \VUgras{umeblowania} \textsuperscript{(87)}.

%%K t11-c2-07-1 p
7. --- \VUgras{Aktorzy} \textsuperscript{(89)} i \VUgras{aktorki}
\textsuperscript{(88)}, \VUgras{śpiewacy} i \VUgras{śpiewaczki} musieli uciekać w swoim
teatralnym \VUgras{kostiumie} \textsuperscript{(90)}. Tenor tłumaczy
\VUgras{dziennikarzowi} \textsuperscript{(92)}, jak to się stało.
\VUgras{Sprawozdawca} przybiegł zaraz od pierwszej chwili z dostojnikami:
\VUgras{prefektem} \textsuperscript{(91)}, \VUgras{burmistrzem} \textsuperscript{(93)},
\VUgras{komisarzem policji} \textsuperscript{(94)} i \VUgras{urzędnikiem sądowym}
\textsuperscript{(95)}, którzy przeprowadzą dochodzenie, żeby osądzić
odpowiedzialności.
\VUsaut{6.0mm}
\VUfilet{22mm}
